\section{A clinical cohort}\label{sec:realdata}

The Diabetes 130-US hospitals cohort \citep{strack2014impact,clore2014uci} records $101{,}766$
admissions of patients with diabetes at $130$ hospitals over ten years, with readmission within
thirty days as the outcome. Patients recur across admissions, so we follow the data set's authors and
keep one admission per patient, the first, excluding admissions that ended in death or discharge to
hospice, whose patients cannot be readmitted, and the three records whose sex is unknown. That leaves
$69{,}987$ patients. We split them in half at random, fitted a logistic model on the first half, and
froze its predictions for the second. The model uses the routinely recorded predictors: age, sex, length
of stay, the numbers of laboratory procedures, other procedures, medications, outpatient, emergency and
inpatient visits in the preceding year and diagnoses, insulin use in four levels, a change of diabetes
medication, and any diabetes medication. No attempt was made to
make the model fit well or badly; it discriminates weakly, with an area under the curve of $0.60$ on
the validation half.

\subsection{The partition on real data}\label{sec:realpartition}

On the development half ($34{,}993$ patients) the fit is checked as the partition is refined from
$G=10$ to $G=2000$ (Figure~\ref{fig:diabetes}(b)). EDGE rejects at every partition,
with $p$ between $4\times10^{-5}$ and $7\times10^{-4}$. The Hosmer--Lemeshow statistic rejects at ten
groups ($p=0.041$) and at no finer partition, with $p$ between $0.10$ and $0.50$ from twenty groups on;
its Farrington correction behaves the same way. Two analysts who chose $G=10$ and $G=50$ would reach
opposite conclusions about the same model \citep{paul2013power}; EDGE's reference does not change
with the number of groups (Theorem~\ref{thm:growingG}).

\subsection{Validation}\label{sec:realvalidation}

On the validation half ($34{,}994$ patients) nothing is refitted, so $\Omega=I$ exactly, and because no
score equation absorbs the intercept, the statistic carries it as a fourth column beside the default
basis. The usual check, the calibration intercept-and-slope test of \citet{cox1958} and
\citet{Miller1991}, rejects: the calibration slope is $0.85$, so the predictions are too spread out, and
the likelihood-ratio statistic is $13.3$ on two degrees of freedom, $p=0.001$. The miscalibration is
small in absolute terms: the integrated calibration index of \citet{austin2019ici} is $0.007$, with E50
$0.007$ and E90 $0.008$, so the predictions are off by less than one percentage point for nine patients in
ten; with $35{,}000$ patients the tests detect it, and whether it matters clinically is a separate
question \citep{nattino2020assessing,vancalster2019}. EDGE
rejects at every partition, with $S$ between $23$ and $41$ on four degrees of freedom and $p$ between
$1\times10^{-4}$ and $3\times10^{-8}$, and the GiViTI belt in its external mode rejects as well
($p=1.5\times10^{-8}$). The Hosmer--Lemeshow statistic rejects at five of the eight partitions, with $p$
from $5\times10^{-6}$ at ten groups to $0.062$, $0.15$ and $0.055$ at five hundred, a thousand and two
thousand groups.

The shape of the miscalibration is not a simple tilt. Figure~\ref{fig:diabetes}(a) shows the
top decile and the second over-predicted and the middle deciles under-predicted, by $1.0$ to $1.5$
percentage points in deciles three, five, six and nine. A grouped test restricted to a tilt, the intercept and slope columns alone, does
not reject at any partition ($p$ between $0.08$ and $0.34$): at ten groups only about $2$ of the
statistic's $23$ lie in the intercept and slope directions, and the rest in the curvature. The grouped
tilt is less sensitive than the record-level Cox--Miller test to a calibration slope of $0.85$, so 
EDGE complements the calibration intercept and slope and does not replace them.

Figure~\ref{fig:belt} draws the ten standardized residuals against the reference the test uses, with
the pointwise band $\pm1.96$ and the simultaneous band $\pm2.80$ over all ten groups. Examining observed
and expected events group by group is standard practice after a Hosmer--Lemeshow test
\citep{hosmer2013applied}; what the figure adds is that the curve drawn through the residuals is the projection $P_Z r$, whose squared length is the statistic, so
the verdict and its location are read from the same picture. Six deciles clear the pointwise band and
one clears the simultaneous one, the sixth, at $r_g=3.1$; the top decile, at $-2.79$, sits on its edge.

For a model used at a threshold the shape has a direct cost. The highest-risk tenth of patients are
predicted to be readmitted at $16.9\%$ and are readmitted at $15.2\%$, $1.7$ percentage points fewer
than predicted, while patients predicted at $7$ to $11\%$ carry a risk up to $1.5$ points higher
than reported. An intercept and slope, the ``weak'' level of the calibration hierarchy
\citep{vancalster2016hierarchy}, record that the model is miscalibrated but not where. A smoothed calibration curve
\citep{vancalster2019,austin2019ici} shows where too, and it remains the first display; EDGE adds a
test whose verdict does not depend on the smoother or the partition, and whose residuals are read on
the same picture.

\begin{figure}[htbp]\centering
\includegraphics[width=\textwidth]{Fig10_diabetes}
\caption{The Diabetes-130 cohort, one admission per patient. (a)~Validation half
($n=34{,}994$), frozen model; observed readmission rate against mean predicted risk in ten
equal-frequency risk groups, with the gap from the diagonal drawn as a bar. (b)~Development
half ($n=34{,}993$), the same model; the $p$-value of EDGE, the Hosmer--Lemeshow statistic
and its Farrington correction as the partition is refined from $G=10$ to $G=2000$, on a
logarithmic scale; the dotted line is $0.05$.}
\label{fig:diabetes}
\end{figure}

\begin{figure}[htbp]\centering
\includegraphics[width=\textwidth]{Fig11_belt}
\caption{The residual band implied by EDGE, on the validation half ($n=34{,}994$, frozen
model, ten equal-frequency risk groups). (a)~The standardized residual
$r_g=(O_g-E_g)/\sqrt{V_g}$ of each decile. Nothing was estimated from these outcomes, so each residual
is standard normal under a calibrated model; the shaded band is the simultaneous $95\%$ band over all
ten groups, a Bonferroni band for ten independent residuals, which is valid because nothing was refitted, and the dashed lines the pointwise band. Triangles mark the deciles beyond the simultaneous
band. The curve is the projection of $r$ onto the test's basis, whose squared length is the statistic.
(b)~The same band on the readmission-rate scale.}
\label{fig:belt}
\end{figure}

\subsection{Corrupted records in this cohort}\label{sec:realcorrupt}

The experiment of Section~\ref{sec:contamination} can be repeated on this cohort. On the real
covariates and the frozen predictions we draw the outcome from the model itself, so that the model is
correct by construction and every rejection is a false alarm. Then $k$ records have one covariate
multiplied by four, with their prediction recomputed at the corrupted value and their outcome kept.
There are $500$ replicates at each of $k=0,10,50,100$ in the validation half; the Monte Carlo standard
error of a rate of $0.05$ is $0.010$.

Corrupting the number of laboratory procedures moves no test at any $k$: at a hundred corrupted
records EDGE reads $0.036$, Stukel's joint score $0.052$ and the cubic calibration test
$0.048$. Its coefficient is small, so a fourfold error moves a record's linear predictor by $0.23$ on
average, less than the spread of the linear predictor across the cohort ($0.37$), and no record becomes
extreme.

The number of previous inpatient admissions is the model's largest coefficient, $0.38$. Only $12\%$ of
patients have any previous admission, so a fourfold error changes about one corrupted record in eight,
but it moves that record's linear predictor by $1.7$ on average, more than four times the cohort's
spread. This is an error that passes a range check, because four previous admissions is a plausible
value. It is also the covariate on which the experiment could show most, and it was chosen for that
reason after the first showed nothing. Of a hundred corrupted records about twelve therefore change, and the separation of
Section~\ref{sec:contamination} appears (Table~\ref{tab:cohortcorrupt}).

\begin{table}[htbp]\centering\small
\caption{False-alarm rate on the cohort's validation half ($n=34{,}994$) when $k$ records have the
number of previous inpatient admissions multiplied by four and the model is otherwise correct by
construction. The grouped tests use frozen predictions in the external mode (four columns); the
default partition has $1400$ groups. $500$ replicates at each $k$.}
\label{tab:cohortcorrupt}
\begin{tabular}{rrrrrr}\hline
$k$ & EDGE (default $G$) & EDGE ($G{=}10$) & HL ($G{=}10$) & Stukel joint & cubic LR \\ \hline
0   & 0.050 & 0.044 & 0.048 & 0.050 & 0.042 \\
10  & 0.042 & 0.042 & 0.050 & 0.102 & 0.080 \\
50  & 0.086 & 0.060 & 0.060 & 0.254 & 0.198 \\
100 & 0.080 & 0.054 & 0.048 & 0.448 & 0.314 \\ \hline
\end{tabular}
\end{table}

About twelve changed records among $34{,}994$ patients raise Stukel's score test to $0.448$ and the
cubic calibration test to $0.314$. At the default partition EDGE rises to $0.086$ with
fifty corrupted records, about six of them changed, and stays below $0.09$; at ten groups it stays
within Monte Carlo error of the nominal level.

A simulation of external validation in one design with five covariates (Supporting Information S7) sets
EDGE against the usual tests there. At ten groups its mean size-adjusted power over ten departures is
$0.415$, against $0.420$ for the GiViTI belt and $0.410$ for Cox's recalibration test, which in most data
sets return the same $p$-value; the two lead by up to $0.14$ on a shift and a threshold term, and EDGE by
up to $0.21$ on a U-shaped term and an interaction. With ten reversed records in $1000$ EDGE reads
$0.082$, below the bar of $0.10$, where Cox's test reads $0.158$ and the belt $0.347$; the
Hosmer--Lemeshow statistic and Spiegelhalter's $z$ stay below it too, with less power ($0.335$ and $0.271$).

\subsection{Monitoring a deployed model}\label{sec:stream}

Once the risk groups are fixed by cut points chosen in advance from reference predictions, the
external-mode statistic depends on the data only through four sums per group: observed events, expected
events, the binomial variance and the count. It can therefore be kept up to date as patients arrive:
each new patient updates one group, and the test is recomputed from the $G$ group summaries without
revisiting earlier records. The streamed statistic equals the external statistic computed in one pass on
all the records with the same cut points, whatever the order and batching of the updates; these fixed
groups drift away from equal size when the risk mix of the monitored patients changes. Provided every
group keeps receiving patients, it keeps its $\chi^2_{d+1}$ reference under a calibrated model
(Supporting Information S8). In a simulation fixed before it was run, at ten groups and with $2000$
replicates, the level was $0.047$ and $0.054$ after $10{,}000$ and $100{,}000$ patients, and $0.050$ and
$0.052$ when every monitored patient came from a risk distribution different from the reference; it
detected an overfitted model, calibration slope $0.8$, in $0.99$ of replicates after $2000$ patients.
The statistic accumulates since deployment, so it answers whether the model has been calibrated over
that whole period; to follow drift, the summaries are restarted each period, a month or a quarter, and
each period is tested on its own.
Adding a hundred patients takes about half a millisecond and the test about $1.5$ milliseconds on one
core, however many patients have been seen. One test at any time is valid, but acting on the first
rejection over repeated looks is not: ten looks at $5\%$ raised the false-alarm rate to $0.21$, and
splitting the level over ten looks fixed in advance, $0.5\%$ each, brought it to $0.030$. In
\texttt{ebrahim.gof} the monitor is \texttt{edge.stream()}.
